\documentclass[11pt]{article}
\usepackage[margin=1in]{geometry}
\usepackage{amsmath,hyperref}
\begin{document}
\begin{center}\large Precursor-emission sensitivity in Stevenson et al. (2022)\end{center}
\noindent\textbf{Reference.} Stevenson et al., \emph{Atmospheric Chemistry and
Physics} 22, 14243--14252.\newline
\href{https://doi.org/10.5194/acp-22-14243-2022}{doi:10.5194/acp-22-14243-2022}.

\paragraph{Method.} The study applies sensitivities from the first Task Force on
Hemispheric Transport of Air Pollution experiment to lockdown-related changes
in NO$_x$, CO, and non-methane volatile organic compound (NMVOC) emissions.
The underlying model experiments perturb 2001 anthropogenic emissions by 20\%.
Eleven models supply regional responses; four supply the global experiments
needed to determine the rest-of-world contribution. Chemical lifetimes are
converted to perturbation lifetimes using each model's methane feedback factor.
The response after one year is obtained from the equilibrium methane response
by multiplying by $1-\exp(-1/\tau_p)$, in which $\tau_p$ is the perturbation
lifetime in years. This represents a one-year forcing, not an instantaneous
change in the methane loss flux. Extra aviation NO$_x$ is added after accounting
for the aviation contribution already present in the scaled experiments.

\paragraph{Findings.} The four-model mean attributes $2.86\pm1.13$ ppb yr$^{-1}$
of the 2020 methane growth-rate increase to precursor-emission reductions,
or $54\pm21$\% of the observed 5.3 ppb yr$^{-1}$ anomaly. NO$_x$ reductions
contribute $4.75\pm1.02$ ppb yr$^{-1}$, offset by CO ($-1.42\pm0.29$) and
NMVOC ($-0.48\pm0.39$) reductions. The additional aviation correction is
0.345 ppb yr$^{-1}$; the total aviation contribution, including what was already
counted, is about 0.93 ppb yr$^{-1}$. Using all available regional models raises
the combined central estimate to 3.16 ppb yr$^{-1}$.

\paragraph{Supplement and uncertainty.} The supplementary workbook provides
lifetimes, methane feedback factors, equilibrium responses, emissions, and
one-year sensitivities. Table S13 reproduces the main attribution table, and its
uncertainty combines species-wise model standard deviations in quadrature.
These errors do not include uncertainty in the lockdown emission estimates or
in transferring sensitivities from 2001 to 2020. Table S11 gives global surface
NO$_x$, CO, and NMVOC changes of $-19.38$ Tg NO$_2$, $-73.38$ Tg CO, and
$-15.65$ Tg C, respectively; aviation NO$_x$ decreases by 0.83 Tg NO$_2$.

\paragraph{Use in this reconstruction.} We extract the workbook's cached numerical
values and retain its species and regional separation. We convert ppb to Tg
using He et al.'s 2.77 Tg ppb$^{-1}$ for a common comparison, explicitly retaining
the distinction between a one-year burden response and $-m_{2019}\Delta k_{\rm OH}$.
The study supplies neither an inferred methane-emission anomaly nor a complete
2019--2022 budget. We therefore compare its chemistry contribution separately
and leave those missing quantities unfilled. Its perturbation lifetime is not
substituted for He's chemical lifetime.
\end{document}
